\documentclass[pdflatex,sn-mathphys-num,Numbered]{sn-jnl}

%%% Standard Packages
\usepackage{amsmath,amssymb,amsfonts}
\usepackage{amsthm}
\usepackage{booktabs}
\usepackage{url}

%%% Theorem environments (SN predefined styles)
%%  thmstyleone:   bold head, italic body (numbered) — theorems
%%  thmstyletwo:   italic head, roman body (numbered) — examples, remarks
%%  thmstylethree: bold head, roman body (numbered)  — definitions
\theoremstyle{thmstyleone}
\newtheorem{theorem}{Theorem}[section]
\newtheorem{lemma}[theorem]{Lemma}
\newtheorem{corollary}[theorem]{Corollary}
\newtheorem{proposition}[theorem]{Proposition}
\newtheorem{conjecture}[theorem]{Conjecture}

\theoremstyle{thmstylethree}
\newtheorem{definition}[theorem]{Definition}

\theoremstyle{thmstyletwo}
\newtheorem{example}[theorem]{Example}
\newtheorem{remark}[theorem]{Remark}

\numberwithin{equation}{section}

%=======================================================================
\begin{document}
%=======================================================================

\title[Exceptional Cases of Sequence A247975]
{Exceptional Cases and Statistical Structure of
Sequence A247975: a Computational Study of
Sun's Conjecture~4.1(i)}

\author*[1]{\fnm{Carlo} \sur{Corti}}\email{carlo.corti@outlook.com}

\affil*[1]{\orgname{Independent researcher},
  \orgaddress{\postcode{21-040}, \city{Kalin\'owka}, \country{Poland}}}

\pacs[MSC Classification]{11A41, 11Y55, 11N05}

\abstract{We present a computational investigation of sequence A247975 in the
OEIS, which encodes Sun's Conjecture~4.1(i): for every positive integer
$m$, there exists a positive integer~$n$ such that $(m+n)$ divides
$p_m^2 + p_n^2$, where $p_k$ denotes the $k$-th prime.
We determine $a(m)$ for $119\,995$ of the values
$m = 1, 2, \ldots, 120\,000$, with final search bound
$n \leq 2 \times 10^{11}$;
the remaining $5$ values are unresolved within this bound.
We prove that the modular filter used in the search eliminates an
asymptotically total proportion of candidates, and establish
computational lower bounds $a(m) > 2 \times 10^{11}$ for the five
unresolved cases.
Among the resolved values we identify $606$
exceptional cases (where $a(m) > 10^8$), including
the outstanding Cram\'er-type ratio $a(m)/(m \log m) \approx 196\,577$
at $m = 23\,995$ and the largest absolute value
$a(82\,300) = 81\,961\,045\,941$.
Statistical analysis yields a Hill statistic
$\hat{\alpha}_{\mathrm{Hill}} \approx 0.713$, consistent with
heavy-tailed behaviour; a density-based heuristic model
fails quantitatively, and we identify three structural sources
of discrepancy.
The five unresolved values are proposed as computational challenge
instances.
Divisibility was independently verified for all $606$ exceptional
cases via \texttt{SymPy}; minimality was independently verified
for $319$ cases via an independent C~sieve.}

\keywords{Prime numbers, Integer sequences, OEIS,
Computational number theory, Sun's conjecture, Cram\'er ratio}

\maketitle

%=======================================================================
\section{Introduction}
\label{sec:intro}
%=======================================================================

In a paper submitted to arXiv in 2013 and published in 2019,
Zhi-Wei Sun~\cite{Sun2019} proposed a collection of conjectures on
addition and divisibility involving primes and other arithmetic
functions.
Among these, Conjecture~4.1(i) (p.~15) states the following.

\begin{conjecture}[Sun, 2013/2019]\label{conj:sun}
	For every positive integer $m$, there exists a positive integer $n$
	such that
	\[
	(m + n) \mid p_m^2 + p_n^2,
	\]
	where $p_k$ denotes the $k$-th prime.
\end{conjecture}

The conjecture was submitted to the OEIS in September 2014 and appears
as sequence A247975~\cite{OEIS_A247975}.
A closely related sequence, A248354~\cite{OEIS_A248354}, encodes the
analogous conjecture for $p_{m^2}^2 + p_{n^2}^2$ (Conjecture~4.1(ii)
in~\cite{Sun2019}).
The sequence $a(m)$, defined as the least positive integer $n$
satisfying the divisibility condition, was computed for $m \leq 10\,000$
by Chai Wah Wu~\cite{WuOEIS}.
To the best of our knowledge, no other published computational
extension of sequence~A247975 beyond Wu's $m \leq 10\,000$ exists in
the literature.

The present work has three main contributions.
First, we extend the computation to $m = 1, \ldots, 120\,000$,
providing the largest verified dataset for this conjecture to date.
Second, we carry out a systematic statistical analysis of the
distribution of $a(m)$.
Third, we identify and characterize $5$ values of $m$ for which the
conjecture resists verification within our final search bound
$n \leq 2 \times 10^{11}$, all exhibiting elevated modular
obstruction rates and extreme Cram\'er-ratio lower bounds;
these are proposed as explicit open problems.

%=======================================================================
\section{Setup and Notation}
\label{sec:setup}
%=======================================================================

We reformulate the divisibility condition in modular arithmetic
and establish the number-theoretic constraints underlying the
computational search.

For positive integers $m$ and $n$, set $s = m + n$.
The divisibility condition $(m+n) \mid p_m^2 + p_n^2$ is equivalent to
\[
p_n^2 \equiv -p_m^2 \pmod{s}.
\]
When $\gcd(p_m, s) = 1$, this simplifies to
\[
\left( p_n p_m^{-1} \right)^2 \equiv -1 \pmod{s},
\]
so (under the assumption $\gcd(p_m, s) = 1$) a solution exists at
index $n$ if and only if $p_n$ falls in one of
the admissible residue classes of the form $r \cdot p_m \pmod{s}$,
where $r$ ranges over the square roots of $-1$ modulo~$s$
and $\omega(s)$ denotes the number of distinct prime factors of~$s$.
When such square roots exist, their number is $2^{\omega(s)}$ for
odd~$s$ and $2^{\omega(s)-1}$ for $s \equiv 2 \pmod{4}$
(see below).

When $\gcd(p_m, s) > 1$, i.e., $p_m \mid s = m+n$, the divisibility
chain $p_m \mid s$, $p_m \mid p_m^2$, and $(m+n) \mid (p_m^2 + p_n^2)$
together force $p_m \mid p_n^2$, hence $p_m \mid p_n$ (since $p_m$ is
prime), hence $p_m = p_n$ (since both are prime), so $n = m$.
In this branch the \emph{only} possible solution at modulus $s$ is
$n = m$, and this case is excluded for $m > 1$ by the following
argument.

If $n = m$ were a solution it would require $2m \mid 2p_m^2$,
i.e., $m \mid p_m^2$.
Since $p_m$ is prime, $m \mid p_m^2$ implies $m \in \{1, p_m, p_m^2\}$.
The case $m = p_m^2$ is impossible for $m \geq 2$: since $p_m \geq 2$
we have $p_m^2 \geq 2p_m > 2m > m$, a contradiction.
The case $m = p_m$ is ruled out by an elementary observation: since
$p_1 = 2, p_2 = 3, \ldots$ is a strictly increasing sequence of
integers starting at~$2$, we have $p_m \geq m + 1 > m$ for all
$m \geq 1$, so $p_m \neq m$.
Hence $n = m$ is not a valid solution for any $m \geq 2$.
For $m = 1$, by contrast, $n = m = 1$ is a valid solution: $m + n = 2$
divides $p_1^2 + p_1^2 = 8$, so $a(1) = 1$.
The algorithm handles this edge case before the main search.
For $m \geq 2$, the $n = m$ branch is excluded and the algorithm
proceeds to the main search loop.

A classical result (see \cite{IrelandRosen}, Ch.~5) gives:

\begin{proposition}\label{prop:qr}
	$-1$ is a quadratic residue modulo $s$ if and only if both of the
	following hold:
	\begin{enumerate}
		\item[(a)] $4 \nmid s$ \textup{(}equivalently, $v_2(s) \leq 1$,
		where $v_2(s)$ denotes the 2-adic valuation of $s$\textup{)}, and
		\item[(b)] no prime $q \equiv 3 \pmod{4}$ divides $s$.
	\end{enumerate}
\end{proposition}

\begin{remark}\label{rem:jacobi}
	Condition~(b) is strictly stronger than requiring the Jacobi symbol
	$\left(\frac{-1}{s}\right) = 1$.
	For odd $s$, the Jacobi symbol $\left(\frac{-1}{s}\right) = 1$ if and
	only if the \emph{sum} $E = \sum_{q \equiv 3 \pmod{4}} v_q(s)$ of the
	exponents of all primes $q \equiv 3 \pmod{4}$ dividing $s$ is even
	(including $E = 0$).
	By contrast, actual solvability of $x^2 \equiv -1 \pmod{s}$ requires
	that no such prime divides $s$ at all, i.e., $E = 0$.
	Thus the Jacobi symbol can equal~$1$ even when $-1$ is \emph{not} a
	quadratic residue modulo~$s$; a simple example is $s = 21 = 3 \times 7$,
	for which $E = 1 + 1 = 2$ is even so $\left(\frac{-1}{21}\right) = 1$,
	yet $-1$ is not a quadratic residue modulo $21$ because condition~(b)
	fails.
	The Jacobi symbol is therefore only a necessary condition for
	solvability, not a sufficient one.
	For $s \equiv 2 \pmod{4}$, the Jacobi symbol is undefined (it is only
	defined for odd positive denominators); solvability of
	$x^2 \equiv -1 \pmod{s}$ in this case requires that no odd prime
	$q \equiv 3 \pmod{4}$ divides $s$, which is exactly condition~(b)
	applied to the odd part of $s$.
\end{remark}

When $-1$ is a QR mod $s$ and $s$ is odd (so every prime divisor of
$s$ is $\equiv 1 \pmod{4}$), there are exactly $2^{\omega(s)}$ square
roots of $-1$ modulo $s$, where $\omega(s)$ denotes the number of
distinct prime factors of~$s$ (see \cite{IrelandRosen}, Ch.~5,
Theorem~9).
When $s \equiv 2 \pmod{4}$ (i.e., $v_2(s) = 1$), condition~(a) is
satisfied and the factor of~$2$ contributes exactly one solution to
$x^2 \equiv -1 \pmod{2}$, so the total count is $2^{\omega(s)-1}$;
this case is handled correctly by the filter
(see Section~\ref{sec:filter}).

The modular filter in our implementation checks condition~(a) exactly
and condition~(b) partially (up to a trial-division bound); see
Section~\ref{sec:filter} for details.
Candidates not eliminated by the filter are tested by direct arithmetic.

%=======================================================================
\section{Computational Methods}
\label{sec:methods}
%=======================================================================

\subsection{Algorithm and modular filter}
\label{sec:filter}

Our implementation (written in C) uses a segmented sieve of
Eratosthenes~\cite{CrandallPomerance} to build a table of all primes up
to a specified limit, together with a modular filter to skip values of
$n$ for which $s = m + n$ admits no square root of $-1$, before testing
the divisibility condition directly.

The filter checks, for each candidate $n$, two conditions:
\begin{enumerate}
	\item[(a)] $4 \nmid s$, i.e., $s \not\equiv 0 \pmod{4}$
	(condition~(a) of Proposition~\ref{prop:qr}, checked exactly); and
	\item[(b$'$)] no prime $q \equiv 3 \pmod{4}$ up to a trial-division
	bound $B = 10^4$ divides $s$ (a partial check of
	condition~(b) of Proposition~\ref{prop:qr}).
\end{enumerate}
This is a safe filter: it eliminates only values of $s$ for which
$x^2 \equiv -1 \pmod{s}$ provably has no solution, so no valid
candidate is ever discarded.
More precisely, if either condition~(a) or~(b) of
Proposition~\ref{prop:qr} is violated for~$s$, then no
$x \in \mathbb{Z}$ satisfies $x^2 \equiv -1 \pmod{s}$, and
consequently the divisibility $(m+n) \mid p_m^2 + p_n^2$ cannot hold
when $\gcd(p_m, s) = 1$ (since it would require
$(p_n p_m^{-1})^2 \equiv -1 \pmod{s}$); the filter is therefore
\emph{sound} in the sense that it discards only provably impossible
candidates.
Note that for $m \geq 2$ the case $\gcd(p_m, s) > 1$ would force
$n = m$, which is not a valid solution (as shown in
Section~\ref{sec:setup}); the filter may therefore assume
$\gcd(p_m, s) = 1$ for all remaining candidates, and its conditions
correctly characterise impossibility.
Values of $s$ that have a prime factor $q \equiv 3 \pmod{4}$ larger
than $B$ are not eliminated by the filter; they pass through to the
direct divisibility test, which may still fail, resulting in a small
amount of redundant arithmetic.
Across the range $m = 1, \ldots, 120\,000$, the filter eliminates
empirically $73$--$77\%$ of candidate values of $n$ before any
arithmetic is performed; we call this per-$m$ elimination percentage
the \emph{obstruction rate}.
The following proposition makes the asymptotic obstruction rate precise.

\begin{proposition}[Asymptotic obstruction rate]\label{prop:obstruction}
Let $\mathcal{A}(x)$ denote the set of integers $s \leq x$ satisfying
both conditions of Proposition~\textup{\ref{prop:qr}}, i.e.,
$v_2(s) \leq 1$ and $s$ has no prime factor $q \equiv 3 \pmod{4}$.
Then
\[
\frac{|\mathcal{A}(x)|}{x}
\;\sim\;
\frac{C_0}{\sqrt{\ln x}}
\qquad (x \to \infty),
\]
where $C_0 > 0$ is an explicit constant.
In particular, the proportion of candidates eliminated by the exact
filter \textup{(}conditions \textup{(a)} and \textup{(b)} of
Proposition~\textup{\ref{prop:qr})} tends to~$1$ as the search
bound grows.
\end{proposition}

\begin{proof}
An integer $s$ satisfies condition~(b) if and only if every prime
factor of $s$ lies in the set $\mathcal{P} = \{2\} \cup
\{p \text{ prime} : p \equiv 1 \pmod{4}\}$.
The classical $B_1$ numbers (integers representable as a sum
of two squares) are characterised by the condition that every prime
$q \equiv 3 \pmod{4}$ appears to an \emph{even} exponent in the
factorisation; our set $\mathcal{A}(x)$ imposes the strictly
stronger requirement that no such prime divides~$s$ at all,
so $\mathcal{A}(x)$ is a proper subset of the $B_1$ set.
By the prime number theorem for arithmetic progressions, the primes
in $\mathcal{P}$ have natural density $\frac{1}{2}$ among all primes.
A theorem of Wirsing~\cite{Wirsing1961} on sums of non-negative
multiplicative functions gives the count of integers up to~$x$ whose
prime factors all belong to a set of primes with
density $\delta$ as $\sim C \cdot x/(\ln x)^{1-\delta}$
for an explicit constant $C > 0$ depending on the set.
Applying this with $\delta = \frac{1}{2}$ yields
$|\{s \leq x : \text{condition (b) holds}\}| \sim C_1 \cdot x/\sqrt{\ln x}$;
the asymptotic order matches the classical Landau--Ramanujan
result~\cite{Landau1908} for the larger $B_1$ set, but the
constant $C_1$ is smaller since fewer integers are counted.
Imposing condition~(a) additionally ($v_2(s) \leq 1$) excludes those
$s$ divisible by~$4$.
Among integers whose prime factors lie in $\mathcal{P}$,
the proportion with $v_2(s) \geq 2$ is asymptotically constant,
so the combined count retains
the same order $\sim C_0 \cdot x/\sqrt{\ln x}$ with a reduced
constant $C_0 \leq C_1$.
The claimed obstruction rate follows:
$1 - |\mathcal{A}(x)|/x \to 1$ as $x \to \infty$.
\end{proof}

% computation from formal proof; defined C_3 explicitly; corrected
% rounding; resolved variable B overloading.
\begin{remark}\label{rem:C0value}
The constant $C_0$ in Proposition~\textup{\ref{prop:obstruction}} can
be estimated numerically via a Mertens-type product.
Define
\[
C_3 \;=\; \lim_{y \to \infty}\;
\sqrt{\ln y}\;
\prod_{\substack{q \leq y \\ q \equiv 3 \pmod{4}}}
\left(1 - \frac{1}{q}\right),
\]
which exists by the analogue of Mertens' third theorem for primes in
arithmetic progressions
\textup{(}see \textup{\cite{Tenenbaum1995}}, Ch.~I.3\textup{)}.
Computing the partial products up to $y = 10^6$ gives
$C_3 \approx 0.8689$, stable to six decimal places.
Combining with the factor $\frac{3}{4}$ from condition~\textup{(a)}
yields $C_0 \approx \frac{3}{4} \times 0.8689 \approx 0.652$.
\end{remark}

For the \emph{truncated} filter used in our implementation,
which checks condition~(b) only up to a trial-division bound
$B = 10^4$, the survival proportion is
\[
\frac{3}{4} \prod_{\substack{q \leq B \\ q \equiv 3 \pmod{4}}}
\left(1 - \frac{1}{q}\right)
\;\approx\; 0.215,
\]
yielding a predicted obstruction rate of $\approx 78.5\%$.
This is consistent with, though slightly higher than, the
observed range $73$--$77\%$; the residual gap reflects
correlations among consecutive values $s = m+n$ arising from
shared small prime factors and the arithmetic structure of
the sequence $\{m+1, m+2, \ldots\}$.
We refer to this per-$m$ elimination percentage as the
\emph{obstruction rate}; see Section~\ref{sec:notfound}
and Table~\ref{tab:notfound} for the obstruction rates of the
unresolved cases.

All integer arithmetic uses 64-bit signed integer types
(\texttt{int64\_t}) throughout the main computation (Phase~6).
The divisibility check $(p_m^2 + p_n^2) \bmod s = 0$ is implemented
using explicit reduction before multiplication:
\[
\bigl((p_m^2 \bmod s) + (p_n \bmod s)^2\bigr) \bmod s \;=\; 0.
\]
The formula treats $p_m$ and $p_n$ asymmetrically for a concrete
reason: for $m \leq 120\,000$ we have $p_m \leq 1\,583\,539$
(the exact value of $p_{120\,000}$),
so $p_m^2 \leq 2.51 \times 10^{12}$ fits directly in \texttt{int64\_t}
and is reduced after squaring; by contrast, $p_n$ can reach
$4.71 \times 10^{10}$ (at $n = 2 \times 10^9$),
so $p_n^2 \approx 2.2 \times 10^{21}$ would
overflow \texttt{int64\_t} unless $p_n$ is reduced \emph{before}
squaring.
After reduction, $\mathtt{pnr} = p_n \bmod s < s \leq m + n \leq
2.0001 \times 10^9$, and hence
\[
\mathtt{pnr}^2 < (2.0001 \times 10^9)^2
= 4.0004 \times 10^{18}
< 2^{63} - 1 \approx 9.22 \times 10^{18}.
\]
To close the overflow argument: the sum
$(p_m^2 \bmod s) + \mathtt{pnr}^2$ satisfies
$(p_m^2 \bmod s) \leq s - 1 < 2.0001 \times 10^9$
and $\mathtt{pnr}^2 \leq 4.0004 \times 10^{18}$,
so the sum is at most $4.0004 \times 10^{18} + 2.0001 \times 10^9
< 2^{63} - 1$, and the final reduction is safe.
No overflow occurs on \texttt{int64\_t}; neither \texttt{\_\_int128}
nor arbitrary-precision arithmetic is required.

In the extended search (Section~\ref{sec:extended}), the search bound
reaches $n \leq 2 \times 10^{11}$, so
$s = m + n$ can reach $\approx 2 \times 10^{11}$ and the
product $p_n^2 \bmod s$ requires intermediate values that may exceed
$2^{63}$.
Accordingly, the extended-search program uses \texttt{\_\_uint128\_t}
for the modular squaring step $p_n^2 \bmod s$ throughout;
peak RAM is kept below $100$\,MB by replacing the global prime
table with a segmented sieve that computes primes on the fly;
see Section~\ref{sec:extended} for details.

Prime generation uses a bitset segmented sieve of Eratosthenes
storing one bit per odd number, which halves memory relative to a
full-integer representation and is fully deterministic: every number
up to the sieve limit is either confirmed prime or confirmed composite;
no probabilistic primality test is employed.

\subsection{Computation phases}
\label{sec:phases}

The computation proceeded in phases:
\begin{enumerate}
	\item \emph{Phases~1--2} ($m \leq 10\,000$, bound $n \leq 10^8$):
	established baseline results, resolving $9998$ of $10\,000$ cases.
	\item \emph{Phase~3} ($m \leq 10\,000$,
	bound $n \leq 7.6 \times 10^8$):
	resolved the remaining $2$ cases, confirming
	$a(4703) = 760\,027\,770$ (previously noted by Sun~\cite{Sun2019})
	and establishing $a(5263) = 420\,423\,462$.
	\item \emph{Phases~4--6} ($m \leq 120\,000$,
	bound $n \leq 2 \times 10^9$):
	extended the dataset to $m = 120\,000$, resolving $119\,931$ of
	$120\,000$ cases, with $69$ values unresolved at the Phase~6 bound.
\end{enumerate}

All computations were performed on a personal workstation equipped with
an AMD Ryzen~9 7940HS processor and 64\,GB of RAM\@.
The final phase (Phase~6, $m \leq 120\,000$) required approximately
$9$~minutes of wall-clock time and approximately $20.8$\,GB of
RAM: $\sim 3.2$\,GB for the packed-bit sieve up to
$5.14 \times 10^{10}$ (one bit per odd number) and $\sim 17.6$\,GB for
the prime index table storing all $2\,176\,824\,767$ primes below
$5.14 \times 10^{10}$ as 8-byte integers.
The search bound $n \leq 2 \times 10^9$ corresponds to
$p_n \leq 47\,055\,833\,459 \approx 4.71 \times 10^{10}$.
The sieve is built to $5.14 \times 10^{10}$, a margin of
approximately $9\%$ above $p_{2 \times 10^9}$, to provide a
comfortable buffer that avoids edge effects and ensures all needed
primes are available.

\subsection{Extended search}
\label{sec:extended}

After the initial computation (Phase~6) left $69$ values unresolved
at bound $n \leq 2 \times 10^9$, a targeted extended search
progressively extended the bound to $n \leq 2 \times 10^{11}$,
resolving $64$ of the $69$ cases and reducing the unresolved count
to~$5$.

The extended search used a new program (\texttt{sun\_ext\_v11.c})
with a fundamentally different architecture from the Phase-6 program:
instead of building a global prime table (which would require
$\sim 20$\,GB of RAM at bound $10^{10}$, and far more at
$2 \times 10^{11}$), it uses a segmented sieve with no persistent
prime table, computing primes on the fly within segments of
$4 \times 10^6$ odd slots ($0.5$\,MB per thread).
The starting prime $p_{\mathrm{START\_N}}$ is computed once; all
subsequent primes are enumerated within segments.
Peak RAM stays below $100$\,MB regardless of the search bound.
Parallelism is achieved via OpenMP on $16$ threads.

The modular squaring step uses a 128-bit integer type
(\texttt{\_\_uint128\_t}) for $p_n^2 \bmod s$ throughout the
extended search (where $p_n^2$ may exceed $2^{63}$);
a single \texttt{DIVQ} inline-assembly instruction handles the
final modular reduction efficiently.

Table~\ref{tab:extended} summarises the extended search.
The search was halted after Session~4 ($n \leq 2 \times 10^{11}$),
at which point two consecutive sessions (S3 and S4) had each resolved
zero additional cases, strongly suggesting that the remaining $5$
unresolved values likely require $n$ well beyond $2 \times 10^{11}$.

\begin{table}[ht]
	\centering
	\caption{Extended search: progressive resolution of the $69$ cases
		unresolved after the initial computation (Phase~6).
		``+Resolved'' denotes newly resolved cases in
		each session.}
	\label{tab:extended}
	\footnotesize
	\begin{tabular}{llrrrl}
		\toprule
		Session & Bound $n$ & +Resolved & Still open & Time (min) & Code \\
		\midrule
		Phase~6 (initial) & $\leq 2 \times 10^9$   & --- & 69 &  8.8 & \texttt{sun\_v6.c} \\
		S1                & $\leq 5 \times 10^{10}$ &  62 &  7 & 49.0 & \texttt{sun\_ext\_v11.c} \\
		S2                & $\leq 10^{11}$          &   2 &  5 & 68.8 & \texttt{sun\_ext\_v11.c} \\
		S3                & $\leq 1.5 \times 10^{11}$ & 0 &  5 & 102.8 & \texttt{sun\_ext\_v11.c} \\
		S4                & $\leq 2 \times 10^{11}$ &   0 &  5 & 135.1 & \texttt{sun\_ext\_v11.c} \\
		\bottomrule
	\end{tabular}
\end{table}

The total wall-clock time for Sessions S1--S4 was approximately
$356$ minutes ($\approx 6$~hours).

\subsection{Use of AI tools}
\label{sec:ai-disclosure}

The computational implementation, including the C~search program
and verification sieve, was developed with extensive assistance
from Anthropic's Claude language models
(\texttt{claude-sonnet-4-6} and \texttt{claude-opus-4-6}).
Claude also assisted substantially in the statistical analysis
and in the preparation of this manuscript.
All computations were executed and independently verified by the
author, who bears full scientific responsibility for all results,
including those portions where AI assistance was used.

\subsection{Independent verification}
\label{sec:verification}

Results were verified by three methods, as described below.
Minimality of $a(m)$ is fully independently verified for $319$
values; for all other values it depends on the correctness of the
exhaustive C~search.

\begin{enumerate}
	\item \emph{Cross-consistency}: all $10\,000$ values for
	$m \leq 10\,000$ agree exactly between Phase~3 and Phase~6 outputs,
	confirming implementation stability and determinism across
	independent runs.
	Note that this method confirms reproducibility of the algorithm,
	not its correctness; two runs of the same implementation will
	agree regardless of any systematic bug.
	\item \emph{Direct divisibility via \texttt{SymPy}}: for each of the
	$606$ exceptional resolved cases ($a(m) > 10^8$), the values $p_m$
	and $p_{a(m)}$ were computed independently using
	\texttt{sympy.prime()}~\cite{SymPy}, which implements its own
	prime-generation algorithm entirely separate from the C~sieve, and
	we verified $(p_m^2 + p_n^2) \bmod (m+n) = 0$ directly.
	All $606$ verified without error: $542$ cases from the initial
	computation and all $64$ cases from the extended search.
	For the $64$ extended-search cases, $p_{a(m)}$ values range from
	$\approx 4.8 \times 10^{10}$ to $\approx 2.2 \times 10^{12}$;
	these were computed via \texttt{sympy.primepi()} combined with a
	walk using \texttt{sympy.nextprime()}/\texttt{prevprime()}, starting
	from a PNT estimate.
	The \texttt{SymPy} and Python scripts take as input only the values
	of $m$ and the claimed $a(m)$; they do not use the prime table
	produced by the C~program.
	This method confirms the divisibility condition for the claimed
	values but \emph{not} minimality; for exceptional cases not covered
	by Method~3 below, minimality depends on the correctness of the
	C~implementation.
	\item \emph{Independent sieve}: an independent C implementation
	(\texttt{verify\_min\_ext.c}) using a segmented sieve of Eratosthenes,
	with no code shared with the search programs, verified $319$ cases
	by exhaustive search from $n = 1$, confirming both the divisibility
	condition \emph{and} the minimality of $a(m)$.
	The $319$ cases comprise six groups ($318$ cases total)
	plus one smoke-test reference case:
	\begin{itemize}
		\item \emph{Group~B1}: the $100$ smallest exceptional cases,
		with $a(m)$ ranging from $1.00 \times 10^8$ to
		$1.52 \times 10^8$, providing the first independent
		minimality check for exceptional values.
		\item \emph{Group~B2}: $150$ randomly selected non-exceptional
		cases ($a(m) < 10^8$, seed~$= 42$).
		\item \emph{Group~C}: all $40$ resolved cases with
		$a(m) > 10^9$ from the initial computation, including
		$m = 14\,740$
		($a(m) = 1{,}994{,}463{,}433$) and $m = 16\,167$
		($a(m) = 1{,}881{,}095{,}843$), the two largest values in
		that dataset.
		\item \emph{Group~E1}: the $4$ extended-search cases with the
		largest Cram\'er ratios
		($m \in \{22\,802,\ 23\,995,\ 53\,963,\ 82\,300\}$,
		$a(m) \in [4.4 \times 10^{10},\; 8.2 \times 10^{10}]$),
		with individual run times of $1\,736$--$2\,813$ seconds.
		\item \emph{Group~E2}: $15$ uniformly sampled extended-search
		cases with $a(m) \in [2.0 \times 10^9,\; 7.9 \times 10^9)$,
		and \emph{Group~E3}: all $9$ extended-search cases with
		$a(m) \in [1.1 \times 10^{10},\; 3.0 \times 10^{10}]$.
		Together these cover the full range of the extended-search
		dataset below the Group~E1 threshold.
	\end{itemize}
	In addition, a smoke-test case ($m = 11\,924$,
	$a(m) = 6\,119\,832\,581$, elapsed $149$~s) was run as a
	reference check before each group E1--E3 campaign.
	All $319$ cases were confirmed without discrepancy.
	The \texttt{verify\_min\_ext.c} program uses \texttt{\_\_uint128\_t}
	arithmetic for the divisibility check to avoid overflow for
	$p_n^2 \sim 10^{25}$ at the largest extended-search values.
\end{enumerate}

In total, $10\,925$ individual verifications were performed across all
three methods: $10\,000$ cross-consistency checks (Method~1),
$606$ SymPy divisibility checks (Method~2), and $319$ independent
sieve minimality checks (Method~3), with zero discrepancies.

%=======================================================================
\section{Results}
\label{sec:results}
%=======================================================================

\subsection{Summary}
\label{sec:summary}

For $m = 1, \ldots, 120\,000$, we determined $a(m)$ in $119\,995$ cases:
$119\,931$ in the initial computation (Phase~6, bound
$n \leq 2 \times 10^9$) and $64$ additional cases in the extended
search (bound $n \leq 2 \times 10^{11}$; see
Section~\ref{sec:extended}).
The remaining $5$ values are listed and discussed in
Section~\ref{sec:notfound}.
Table~\ref{tab:top10} lists the ten largest values of $a(m)$ found.

We call a resolved $m$ \emph{exceptional} if $a(m) > 10^8$.
This is an author-defined cutoff chosen for exposition;
it is approximately $1300$ times the median value of $a(m)$
($\approx 79\,000$, computed from the $119\,995$ resolved values).
There are $E_{\mathrm{res}}(1.2 \times 10^5) = 606$ such resolved exceptional
cases ($542$ from the initial computation plus all $64$ cases resolved
in the extended search, each of which has $a(m) > 2 \times 10^9$).
Since the $5$ unresolved values all satisfy
$a(m) > 2 \times 10^{11}$ by definition of the final search bound,
the true exceptional count satisfies $E(1.2 \times 10^5) \geq 611$, where
\[
E(M) \;=\; E_{\mathrm{res}}(M)
+ \#\bigl\{m \leq M : m \text{ unresolved}\bigr\}
\]
denotes the total exceptional count; the quantity $E_{\mathrm{res}}(M)$
counts only resolved exceptional cases throughout this paper.

The complete dataset (all $119\,995$ resolved values and the list of
$5$ unresolved~$m$) is publicly available; see
Section~\ref{sec:data}.

\begin{table}[ht]
	\centering
	\caption{Ten largest values of $a(m)$ for $m \leq 120\,000$.
		The Cram\'er-type ratio is $C(m) = a(m)/(m \log m)$
		(natural logarithm throughout).
		All entries were found in the extended search
		(Section~\ref{sec:extended}).}
	\label{tab:top10}
	\begin{tabular}{rrrr}
		\toprule
		$m$ & $a(m)$ & $p_m$ & $a(m)/(m \log m)$ \\
		\midrule
		82\,300 & 81\,961\,045\,941 & 1\,052\,459 & 87\,990 \\
		53\,963 & 50\,362\,759\,647 &   665\,179 & 85\,653 \\
		23\,995 & 47\,572\,503\,750 &   274\,453 & 196\,577 \\
		22\,802 & 44\,000\,698\,448 &   259\,577 & 192\,303 \\
		37\,253 & 29\,834\,694\,588 &   443\,873 & 76\,088 \\
		48\,834 & 23\,630\,637\,176 &   596\,419 & 44\,821 \\
		37\,949 & 23\,612\,836\,020 &   453\,181 & 59\,012 \\
		43\,297 & 20\,932\,270\,152 &   523\,007 & 45\,285 \\
		69\,546 & 17\,485\,936\,480 &   876\,191 & 22\,550 \\
		109\,857 & 14\,929\,726\,016 & 1\,438\,973 & 11\,709 \\
		\bottomrule
	\end{tabular}
\end{table}

\subsection{Outstanding cases}
\label{sec:outstanding}

We list the five largest Cram\'er-type ratios
$C(m) = a(m)/(m\log m)$ among resolved $m \leq 120\,000$
in Table~\ref{tab:ratio}.
Since $p_m \sim m\log m$ by the prime number theorem
(see \cite{Davenport}, Ch.~18), the ratio~$C(m)$
is approximately $a(m)/p_m$,
normalising the solution index by the size of the prime~$p_m$.
Under Cram\'er's probabilistic model for
primes~\cite{Cramer1936}, treating primes as independent random
events with local density $1/\log n$ near~$n$, the expected index
at which one encounters a prime satisfying a given congruence
condition of modulus of order~$m$ is of order~$m$; hence $C(m)$
measures how many times larger the actual solution is compared to
this heuristic baseline.
The extended search dramatically reshuffled the ranking:
$m = 23\,995$ now holds the global resolved maximum
($C \approx 196\,577$), far exceeding the previous record holder
$m = 4703$ ($C \approx 19\,111$).
All five entries in the new top~5 were resolved in the extended search.

\begin{table}[ht]
	\centering
	\caption{Top five resolved values by Cram\'er-type ratio
		$a(m)/(m\log m)$.  All entries were found in the extended search.}
	\label{tab:ratio}
	\begin{tabular}{rrrrr}
		\toprule
		Rank & $m$ & $a(m)$ & $p_m$ & $a(m)/(m\log m)$ \\
		\midrule
		1 & 23\,995 & 47\,572\,503\,750 &  274\,453 & 196\,577 \\
		2 & 22\,802 & 44\,000\,698\,448 &  259\,577 & 192\,303 \\
		3 & 82\,300 & 81\,961\,045\,941 & 1\,052\,459 & 87\,990 \\
		4 & 53\,963 & 50\,362\,759\,647 &  665\,179 & 85\,653 \\
		5 & 37\,253 & 29\,834\,694\,588 &  443\,873 & 76\,088 \\
		\bottomrule
	\end{tabular}
\end{table}

For $m = 23\,995$: $p_m = 274\,453$,
$a(m) = 47\,572\,503\,750$,
$s = m + a(m) = 47\,572\,527\,745$.
The divisibility and minimality of this value were independently
confirmed by \texttt{verify\_min\_ext.c} (Group~E1, elapsed
$1\,854$~seconds).
For $m = 11\,924$ (previous resolved record, now outside the top~5):
$p_m = 127\,289$, $a(m) = 6\,119\,832\,581$,
$s = m + a(m) = 6\,119\,844\,505$,
$C \approx 54\,679$; independently verified by both methods.

\subsection{Growth of resolved exceptional cases}
\label{sec:growth}

Let $E_{\mathrm{res}}(M)$ denote the number of resolved values
$m \leq M$ with $a(m) > 10^8$.
Table~\ref{tab:growth} shows the growth of $E_{\mathrm{res}}(M)$ over
successive ranges of $10\,000$.

\begin{table}[ht]
	\centering
	\caption{Count of resolved exceptional cases $a(m) > 10^8$
		by range of $m$ (including extended-search results).}
	\label{tab:growth}
	\begin{tabular}{lr}
		\toprule
		Range of $m$ & Count \\
		\midrule
		$1$--$10\,000$         &   2 \\
		$10\,001$--$20\,000$   &  30 \\
		$20\,001$--$30\,000$   &  25 \\
		$30\,001$--$40\,000$   &  30 \\
		$40\,001$--$50\,000$   &  44 \\
		$50\,001$--$60\,000$   &  52 \\
		$60\,001$--$70\,000$   &  50 \\
		$70\,001$--$80\,000$   &  63 \\
		$80\,001$--$90\,000$   &  67 \\
		$90\,001$--$100\,000$  &  80 \\
		$100\,001$--$110\,000$ &  81 \\
		$110\,001$--$120\,000$ &  82 \\
		\midrule
		Total & 606 \\
		\bottomrule
	\end{tabular}
\end{table}

%=======================================================================
\section{Statistical Analysis}
\label{sec:statistics}
%=======================================================================

\subsection{Heavy tail}
\label{sec:heavytail}

The distribution of $a(m)$ for resolved $m = 1, \ldots, 120\,000$ is
strongly heavy-tailed.
We apply the Hill estimator~\cite{Hill1975} as an exploratory
diagnostic tool; since the data are deterministic and not
i.i.d., this should be understood as descriptive rather than as
a formal inference about an underlying distribution
(see \cite{Clauset2009} for a discussion of the limitations of
power-law fitting in empirical data).
Given $N_s$ order statistics $X_{(1)} \leq \cdots \leq X_{(N_s)}$, the
Hill estimator with threshold parameter $k$ is defined as
\[
\hat{\alpha}_{\mathrm{Hill}}(k)
= \left(
\frac{1}{k} \sum_{i=1}^{k}
\log \frac{X_{(N_s-i+1)}}{X_{(N_s-k)}}
\right)^{-1}.
\]
The Hill statistic, computed over the top $10\%$ of the $119\,931$
values from the initial computation
($k = 11\,993$, threshold $a(m) > 1\,364\,314$), takes the value
\[
\hat{\alpha}_{\mathrm{Hill}} = 0.713.
\]
Standard confidence intervals for the Hill estimator assume i.i.d.\
data and do not apply to this deterministic dataset; the value above
should be interpreted as a descriptive summary only.
The $64$ values from the extended search are all in the extreme tail
($a(m) > 2 \times 10^9$) and represent $0.05\%$ of the dataset;
their inclusion does not materially change the Hill estimate
(the statistic is driven by the top $10\%$ of the initial dataset,
and $64$ additional points in the far tail shift $\hat{\alpha}$ by
less than $0.002$).
The statistic varies with threshold:
\[
\hat{\alpha}(1\%) = 0.817, \quad
\hat{\alpha}(5\%) = 0.732, \quad
\hat{\alpha}(10\%) = 0.713, \quad
\hat{\alpha}(20\%) = 0.676.
\]
The $10\%$ threshold is adopted as the primary estimate,
balancing sensitivity to the tail against stability of the statistic.
The threshold-dependence indicates that the tail is not a pure power
law; the values of $\hat{\alpha}_{\mathrm{Hill}}$ should be interpreted
only as an indication of heavy-tailed behaviour, not as evidence of a
specific parametric form.

\subsection{Cram\'er-type ratio}
\label{sec:cramerratio}

The Cram\'er-type ratio $C(m) = a(m)/(m \log m)$, as defined in
Section~\ref{sec:outstanding}, has median over resolved
$m = 2, \ldots, 120\,000$
(computed from the $119\,995$ resolved values)
\[
\widetilde{C} \approx 0.154,
\]
stable and well-defined across the full range; the $64$ newly resolved
values lie far in the upper tail and do not affect the median.

%=======================================================================
\section{Heuristic Model and Comparison with Data}
\label{sec:heuristic}
%=======================================================================

\subsection{Density argument}
\label{sec:density}

We present below a density-based heuristic for the expected scale
of $a(m)$; as the data will show, the model fails quantitatively
to predict the observed scale, indicating the need for a refined
understanding.

Fix $m$ and consider the search for a valid $n$.
For a given candidate $s = m + n$, a solution exists if and only if
$p_n$ falls in one of the $2^{\omega(s)}$ residue classes of the form
$r \cdot p_m \pmod{s}$, where $r$ ranges over the square roots of
$-1 \pmod{s}$.
By Dirichlet's theorem~\cite{Davenport}, the asymptotic density
of primes in each such class is $\varphi(s)^{-1}$ for fixed $s$,
where $\varphi(s)$ denotes Euler's totient function (the number of
integers in $\{1, \ldots, s\}$ coprime to $s$); we note
that applying this for moduli $s$ that grow up to $\sim 2 \times 10^9$
while checking primes up to $\sim 4.7 \times 10^{10}$ goes beyond the
range of standard uniformity results (Siegel--Walfisz,
see \cite{Davenport}, Ch.~21), and the following heuristic should be
treated as qualitative.
The total proportion of admissible residue classes modulo $s$ is
\[
D(s) \;=\; \frac{2^{\omega(s)}}{\varphi(s)}
\]
for odd~$s$; for $s \equiv 2 \pmod{4}$ the numerator is
$2^{\omega(s)-1}$, a distinction absorbed into the leading constant
in the estimates below.
Note that $D(s)$ is the \emph{proportion of admissible residue classes},
not a probability of success for a randomly chosen prime near $s$;
the following derivation should be regarded as a rough heuristic
envelope rather than a rigorous estimate.

Using the normal-order estimate $\omega(s) \approx \log\log s$
\cite{HardyRamanujan} and the extremal lower bound
$\varphi(s) \gtrsim s \cdot e^{-\gamma}/\log\log s$
\cite{Mertens}
(the first is a normal-order result valid for almost all $s$, the
second an extremal bound realised at primorials; their combination
has no rigorous standing), one obtains:
\[
D(s) \;\lesssim\;
\frac{(\log s)^{\log 2} \cdot e^{\gamma} \log\log s}{s}.
\]
Treating each candidate $n$ as independently selected
(in the spirit of the Cram\'er model, which treats primes as random
with local density $1/\log n$)
with probability $D(m+n)$ of yielding a valid solution,
the expected value of $a(m)$
satisfies $\sum_{n=1}^{N} D(m+n) \approx 1$, which for slowly varying
$D$ reduces to $\int_m^{m+N} D(s)\,ds \approx 1$.
Substituting the bound on $D(s)$ and integrating yields
\[
\int_m^{m+N} \frac{(\log s)^{\log 2} \cdot e^{\gamma} \log\log s}{s}
\,ds
\;\approx\;
e^{\gamma} \cdot (\log m)^{\log 2} \cdot \log\log m
\cdot \frac{N}{m}
\;\approx\; 1,
\]
where the approximation treats $\log s \approx \log m$ and
$\log\log s \approx \log\log m$ over the range $[m, m+N]$, valid
since $N \ll m$ in the regime of interest.
This gives $N \sim m/(e^{\gamma} (\log m)^{\log 2} \log\log m)$;
for $m = 10^5$ this yields
$\approx 0.04\,m$ --- well below the observed median
$a(m) \approx 1.3\,m$.
This large discrepancy shows that the density model does not explain
the observed scale of $a(m)$.
This prediction is purely illustrative and should not be interpreted
as a rigorous asymptotic estimate; the underlying derivation combines
a normal-order result with an extremal bound in a mathematically
unjustified way, as noted above.
The Cram\'er-type normalization
$m\log m$ is motivated independently by the prime-index analogy
(Section~\ref{sec:cramerratio}), not derived from $D(s)$.

\subsection{Comparison with data}
\label{sec:comparison}

The density model provides a qualitative explanation for the heavy tail
(the factor $(\log s)^{\log 2}$ in $D(s)$ produces slower-than-power-law
decay), but fails quantitatively: it predicts the typical scale of
$a(m)$ as $\sim m/(e^{\gamma} (\log m)^{\log 2} \log\log m)
\approx 0.04\,m$, whereas the typical $a(m)$ scales as
approximately $1.3\,m$.
The growing count $E_{\mathrm{res}}(M)$ in Table~\ref{tab:growth} is
qualitatively consistent with increasing exceptional-case frequency,
but the growth factor of $303$ from $E_{\mathrm{res}}(10^4) = 2$ to
$E_{\mathrm{res}}(1.2 \times 10^5) = 606$ is far larger than any simple
log-power formula would predict.
We therefore regard the density model as providing only a qualitative
motivation for heavy-tailed behaviour, and do not propose a
quantitative formula for $E_{\mathrm{res}}(M)$.

\subsection{Sources of discrepancy}
\label{sec:discrepancy}

The density model of Section~\ref{sec:density} underestimates the
typical value of $a(m)$ by a factor of approximately $30$
(predicting $\sim 0.04\,m$ versus observed $\sim 1.3\,m$).
We identify three structural sources of this discrepancy, none of
which is captured by the independence assumptions underlying the
model.

\paragraph{Correlated candidates.}
The model treats each candidate $n$ as an independent trial.
In reality, the moduli $s = m + n$ and $s' = m + n + 1$ are
consecutive integers, and their divisibility by small primes is
strongly correlated: if $q \mid s$ then $q \nmid s'$.
Over a window of $q$ consecutive values of $s$, exactly one is
divisible by~$q$; this short-range anticorrelation means that the
``success events'' (finding a prime in an admissible residue class)
are not independent, and the waiting-time distribution differs from
the geometric distribution assumed by the model.

\paragraph{Varying modulus.}
For each candidate $n$, the modulus $s = m + n$ changes, so the set
of admissible residue classes changes with~$n$.
The density model averages over $s$ as though the residue class
condition were fixed, but in practice the target moves with every
step.
In particular, $D(s)$ fluctuates between zero (when $s$ has a prime
factor $q \equiv 3 \pmod{4}$) and $\gg 1/s$ (when $s$ is a product
of small primes $\equiv 1 \pmod{4}$); this intermittency produces
longer runs of consecutive failures than the averaged model predicts.

\paragraph{Interaction with $p_m$.}
The density $D(s)$ accounts for the proportion of admissible residue
classes modulo~$s$, but the actual condition also requires $p_n$ to
be a prime landing in one of these classes.
The distribution of primes in residue classes modulo~$s$ is governed
by Dirichlet's theorem, which gives equidistribution only
asymptotically; for the finite ranges involved
($p_n \leq 4.7 \times 10^{10}$ at bound $n = 2 \times 10^9$,
while $s$ can be as large as $\sim 2 \times 10^9$), there can be
significant deviations from equidistribution, particularly for
moduli $s$ with special arithmetic properties relative to $p_m$.

These three effects compound to produce a systematic underestimate
of the waiting time.
Collectively, they plausibly explain much of the observed factor of
approximately~$30$ between the heuristic prediction
($\sim 0.04\,m$) and the data ($\sim 1.3\,m$).
A refined model incorporating even the first correction (correlated
candidates via an inclusion-exclusion over small primes) would
likely close much of the gap; see also \cite{Granville1995} for a
discussion of corrections to the classical Cram\'er heuristic
in related contexts.
We leave the development of such a model as a direction for
future work.

The outstanding Cram\'er-type ratio of $m = 23\,995$
(ratio $\approx 196\,577$, the global resolved maximum over
$m \leq 120\,000$) may reflect an especially unfavourable interaction
between $p_{23\,995} = 274\,453$ and the moduli $s = m+n$, producing
a suppression of $D(s)$ not captured by the average model.
Identifying the specific arithmetic mechanism responsible for such
extreme outliers remains an open question.

%=======================================================================
\section{Unresolved Cases}
\label{sec:notfound}
%=======================================================================

The $69$ values of $m \leq 120\,000$ that were unresolved in the
initial computation ($n \leq 2 \times 10^9$) were subjected to the
extended search described in Section~\ref{sec:extended}.
Of these, $64$ were resolved, leaving $5$ values unresolved
within the final search bound $n \leq 2 \times 10^{11}$.
They are listed in Table~\ref{tab:notfound}, together with $p_m$,
the filter-detected obstruction rate, and estimated lower bounds on
the Cram\'er ratio $C(m)$.
The \emph{obstruction rate} is the percentage of candidate values of $n$
(within the search range) eliminated by the modular filter of
Section~\ref{sec:filter} before any arithmetic is performed.

\begin{table}[ht]
\centering
\caption{The $5$ unresolved values of $m \leq 120\,000$, with $p_m$,
	filter-detected obstruction rate (percentage of candidates
	eliminated by the modular filter; see Section~\ref{sec:filter}),
	and estimated Cram\'er ratio lower bound.
	No solution was found with $n \leq 2 \times 10^{11}$.}
\label{tab:notfound}
\small
\begin{tabular}{rrrr}
	\toprule
	$m$ & $p_m$ & Obstr.\ (\%) & $C(m) >$ \\
	\midrule
	37\,249  &   443\,837 & 76.2 & 510\,126 \\
	66\,257  &   831\,301 & 76.9 & 271\,910 \\
	76\,868  &   977\,087 & 77.0 & 231\,280 \\
	98\,379  & 1\,276\,747 & 77.2 & 176\,831 \\
	117\,862 & 1\,552\,879 & 77.4 & 145\,316 \\
	\bottomrule
\end{tabular}
\end{table}

All $5$ cases exhibit obstruction rates in the range $76.2$--$77.4\%$,
consistent with the $75.7$--$77.4\%$ range observed for the original
$69$ unresolved cases.
The extended search resolved $64$ of those $69$ cases
(Table~\ref{tab:extended}), demonstrating that the obstruction rate
alone does not determine whether a case is resolvable --- rather, it is
the Cram\'er ratio that correlates with difficulty.
The $5$ remaining cases are precisely those with the highest estimated
Cram\'er ratios among the original~$69$, all exceeding $145\,000$.

\paragraph{Resolution curve.}
The extended search resolved cases at a strongly decreasing rate:
$51$ cases in the first range $n \in [2 \times 10^9, 10^{10}]$,
$9$ more in $[10^{10}, 3 \times 10^{10}]$,
$2$ in $[5 \times 10^{10}, 10^{11}]$,
and $0$ in the final range $[10^{11}, 2 \times 10^{11}]$.
Two consecutive sessions (S3 and S4) resolved no cases.

\paragraph{Computational lower bounds.}
The exhaustive search establishes the following.

\begin{theorem}[Computational lower bounds]\label{thm:lower}
For each $m \in \{37\,249,\; 66\,257,\; 76\,868,\; 98\,379,\;
117\,862\}$, the divisibility condition
$(m+n) \mid p_m^2 + p_n^2$ has no solution with
$1 \leq n \leq 2 \times 10^{11}$.
Consequently, either $a(m) > 2 \times 10^{11}$ or
Conjecture~\textup{\ref{conj:sun}} fails for this value of $m$.
\end{theorem}

\begin{proof}
The search program \texttt{sun\_ext\_v11.c} tests, for each of the
five values of $m$, every integer $n$ in the range
$1 \leq n \leq 2 \times 10^{11}$: it enumerates primes $p_n$ via a
segmented sieve of Eratosthenes and, for each $n$ not eliminated by
the modular filter of Section~\ref{sec:filter}, evaluates the
divisibility condition $(p_m^2 + p_n^2) \bmod (m+n) = 0$ directly
using $128$-bit arithmetic (Section~\ref{sec:extended}).
The modular filter is \emph{sound}
(Proposition~\ref{prop:qr}): it eliminates only values of $n$ for
which no solution can exist, so no valid candidate is skipped.
For all five values of $m$, the search returns no solution in the
stated range; the sieve is deterministic and complete, so the
conclusion follows.
The source code is publicly available~\cite{GithubRepo}.
\end{proof}

\begin{remark}
The result of Theorem~\textup{\ref{thm:lower}} is conditional on the
correctness of the deterministic search implementation.
The implementation has been independently spot-checked for $319$
values via Method~3 of Section~\textup{\ref{sec:verification}},
with zero discrepancies; this includes the $4$ resolved cases with
the largest Cram\'er ratios.
\end{remark}

In particular, the Cram\'er-type ratio for each of these values
satisfies $C(m) = a(m)/(m \log m) > 2 \times 10^{11}/(m \log m)$,
yielding the lower bounds in Table~\ref{tab:notfound}.

\paragraph{Computational challenge.}
We propose the five values of Theorem~\ref{thm:lower} as
computational challenge instances: for each, determine $a(m)$ or
establish that no solution exists.
The decreasing resolution rate (Table~\ref{tab:extended}) and the
absence of any resolution in the range
$n \in [10^{11}, 2 \times 10^{11}]$ indicate that resolving these
cases will require searching substantially beyond
$n = 2 \times 10^{11}$.
Whether any of these values constitutes a genuine counterexample to
Conjecture~\ref{conj:sun} remains an open question; the pattern of
progressive resolution observed in this study (from $69$ unresolved
at bound $2 \times 10^9$ to $5$ at bound $2 \times 10^{11}$) is
consistent with these being extreme outliers of the tail distribution
rather than structural failures of the conjecture.

%=======================================================================
\section{Conclusion}
\label{sec:conclusion}
%=======================================================================

We have determined $a(m)$ for $119\,995$ values of $m \leq 120\,000$,
with $5$ values remaining unresolved within our final search bound
$n \leq 2 \times 10^{11}$.
The initial computation ($n \leq 2 \times 10^9$) resolved $119\,931$
cases; a targeted extended search resolved $64$ of the
$69$ initially unresolved cases, reducing the count to~$5$.
The resolved dataset reveals a heavy-tailed distribution with Hill
statistic $\hat{\alpha}_{\mathrm{Hill}} \approx 0.713$, a growing
density of exceptional cases, and a new outstanding resolved case at
$m = 23\,995$ whose Cram\'er-type ratio ($\approx 196\,577$) is
the global resolved maximum over the entire range studied, more than
ten times the previous record holder $m = 4703$ ($\approx 19\,111$).

We proved that the modular filter eliminates an asymptotically
increasing proportion of candidates
(Proposition~\ref{prop:obstruction}), with the theoretical rate
for the truncated filter ($\approx 78.5\%$) consistent with the
observed range ($73$--$77\%$).
We identified three structural sources for the quantitative failure
of the density-based heuristic (Section~\ref{sec:discrepancy}):
correlated candidates, varying moduli, and finite-range deviations
from equidistribution in arithmetic progressions.

For the $5$ unresolved cases, we established computational lower bounds
$a(m) > 2 \times 10^{11}$ (Theorem~\ref{thm:lower}) and proposed
them as computational challenge instances.
The progressive resolution of the original $69$ unresolved cases
--- from $69$ at bound $2 \times 10^9$ to $5$ at bound
$2 \times 10^{11}$ --- demonstrates that the conjecture was verified
computationally for all but the most arithmetically resistant values
of $m$ in this range.

The heuristic density model provides a qualitative explanation for the
heavy tail but fails quantitatively to predict the scale of $a(m)$ or
the growth rate of exceptional cases.

\paragraph{Comparison with Conjecture~4.3(i).}
A companion computation~\cite{CortiA248044} applies the same
methodology to sequence A248044~\cite{OEIS_A248044}, which encodes
Sun's Conjecture~4.3(i): $(m+n) \mid \pi(m)^2 + \pi(n)^2$.
Computing $a(m)$ for $m \leq 100\,000$ with bound $n \leq 10^{13}$
resolves $99\,998$ of $100\,000$ cases, leaving only $2$ unresolved.
The comparison reveals instructive structural differences:
the obstruction rate is lower (${\sim}\,57\%$ vs.\ $73$--$77\%$
here), the maximum resolved Cram\'er-type ratio is much larger
(${\sim}\,6.5 \times 10^7$ vs.\ ${\sim}\,1.97 \times 10^5$),
and the two unresolved cases form an arithmetically structured
cluster sharing $\pi(m) = 2225 = 5^2 \times 89$, whose
factorisation into primes $\equiv 1 \pmod{4}$ produces maximal
modular obstruction --- a phenomenon absent in the present sequence,
where $p_m$ is always prime and the obstruction structure is more
uniform.
These contrasts indicate that the difficulty of Sun's divisibility
conjectures depends critically on the arithmetic function involved
($p_k$ vs.\ $\pi(k)$), not only on the shared divisibility
framework $(m+n) \mid f(m)^2 + f(n)^2$.

\paragraph{Future work.}
Natural directions include an extended computational
search for the $5$ remaining unresolved values of A247975
(and the $2$ of A248044), a refined probabilistic
model incorporating the correlation structure identified in
Section~\ref{sec:discrepancy},
and a systematic comparative study across the full family of
Sun's divisibility conjectures~\cite{Sun2019}.

%=======================================================================
\backmatter
%=======================================================================

\bmhead{Acknowledgements}

The author thanks Zhi-Wei Sun for the beautiful conjecture that
motivated this work, and Chai Wah Wu, whose extension of the OEIS
b-file for A247975 to $m \leq 10\,000$ provided a valuable
cross-check for our results.

\section*{Declarations}

\begin{description}
\item[Funding] No funding was received for conducting this study.

\item[Competing interests] The author has no relevant financial
  or non-financial interests to disclose.

\item[Data availability] \label{sec:data}
  The complete dataset ($119\,995$ resolved values and $5$ unresolved
  $m$-values) is publicly available at:
  \begin{itemize}
    \item GitHub repository:
          \url{https://github.com/carcorti/A247975}
    \item Permanent Zenodo archive (DOI):\\
          \url{https://doi.org/10.5281/zenodo.18920371}
    \item OEIS: \url{https://oeis.org/A247975}
  \end{itemize}
  The repository will be synchronized with the submitted version of the
  manuscript and all associated data files prior to publication.

\item[Code availability]
  The C~source code for the main search
  (\texttt{sun\_v6.c}), the extended-search program
  (\texttt{sun\_ext\_v11.c}), the independent verification sieves
  (\texttt{verify\_min.c} and \texttt{verify\_min\_ext.c}), and the
  Python/SymPy verification scripts
  (\texttt{verify\_ext\_sympy.py}, \texttt{run\_verify\_ext\_batch.py})
  are available at the GitHub and Zenodo repositories listed above.

\item[Author contribution] Not applicable (single author).
\end{description}

\begin{thebibliography}{99}

	\bibitem{Sun2019}
	Sun, Z.-W.:
	Some new problems in additive combinatorics.
	Nanjing Univ.\ J.\ Math.\ Biquarterly \textbf{36}(2), 134--155
	(2019).
	Preprint: arXiv:1309.1679 [math.NT] (submitted 2013, v9: 2020).
	Conjecture~4.1(i) appears on p.~15; Remark~4.1 notes the case
	$m = 4703$.
	\url{https://arxiv.org/abs/1309.1679}

	\bibitem{OEIS_A247975}
	OEIS Foundation:
	Sequence A247975.
	The On-Line Encyclopedia of Integer Sequences
	(accessed March~2026).
	\url{https://oeis.org/A247975}

	\bibitem{OEIS_A248354}
	OEIS Foundation:
	Sequence A248354.
	The On-Line Encyclopedia of Integer Sequences
	(accessed March~2026).
	\url{https://oeis.org/A248354}

	\bibitem{Hill1975}
	Hill, B.M.:
	A simple general approach to inference about the tail of a
	distribution.
	Ann.\ Statist.\ \textbf{3}, 1163--1174 (1975)

	\bibitem{IrelandRosen}
	Ireland, K., Rosen, M.:
	A Classical Introduction to Modern Number Theory,
	2nd edn. Springer, New York (1990).
	Ch.~5, Theorem~9 for quadratic residues modulo composite integers

	\bibitem{Davenport}
	Davenport, H.:
	Multiplicative Number Theory,
	3rd edn. Springer, New York (2000).
	Ch.~4 for Dirichlet's theorem; Ch.~21 for the Siegel--Walfisz
	theorem and uniformity in arithmetic progressions

	\bibitem{SymPy}
	Meurer, A., et al.:
	SymPy: symbolic computing in Python.
	PeerJ Computer Science \textbf{3}, e103 (2017).
	Version~1.12 used for independent prime computation

	\bibitem{HardyRamanujan}
	Hardy, G.H., Ramanujan, S.:
	The normal number of prime factors of a number~$n$.
	Quart.\ J.\ Math.\ \textbf{48}, 76--92 (1917)

	\bibitem{Mertens}
	Mertens, F.:
	Ein Beitrag zur analytischen Zahlentheorie.
	J.\ Reine Angew.\ Math.\ \textbf{78}, 46--62 (1874).
	See also \cite{Tenenbaum1995}, Ch.~I.3

	\bibitem{Tenenbaum1995}
	Tenenbaum, G.:
	Introduction to Analytic and Probabilistic Number Theory.
	Cambridge Univ.\ Press (1995)

	\bibitem{Cramer1936}
	Cram\'er, H.:
	On the order of magnitude of the difference between consecutive
	prime numbers.
	Acta Arith.\ \textbf{2}, 23--46 (1936)

	\bibitem{WuOEIS}
	Wu, C.W.:
	Table of $n$, $a(n)$ for $n = 1..10\,000$,
	contribution to OEIS sequence A247975, 2014.
	\url{https://oeis.org/A247975/b247975.txt}

	\bibitem{Clauset2009}
	Clauset, A., Shalizi, C.R., Newman, M.E.J.:
	Power-law distributions in empirical data.
	SIAM Review \textbf{51}, 661--703 (2009)

	\bibitem{CrandallPomerance}
	Crandall, R., Pomerance, C.:
	Prime Numbers: A Computational Perspective,
	2nd edn. Springer, New York (2005).
	Ch.~3 for sieve algorithms

	\bibitem{Landau1908}
	Landau, E.:
	\"Uber die Einteilung der positiven ganzen Zahlen in vier
	Klassen nach der Mindestzahl der zu ihrer additiven Zusammensetzung
	erforderlichen Quadrate.
	Arch.\ Math.\ Phys.\ \textbf{13}, 305--312 (1908).
	Establishes the asymptotic count ${\sim}\,Kx/\sqrt{\ln x}$
	(with $K \approx 0.7642$) of integers up to $x$ with no prime factor
	$q \equiv 3 \pmod{4}$ to odd power; the natural density of this
	set is zero

	\bibitem{Granville1995}
	Granville, A.:
	Harald Cram\'er and the distribution of prime numbers.
	Scand.\ Actuarial J.\ \textbf{1995}(1), 12--28 (1995)

	\bibitem{Wirsing1961}
	Wirsing, E.:
	Das asymptotische Verhalten von Summen \"uber multiplikative
	Funktionen.
	Math.\ Ann.\ \textbf{143}, 75--102 (1961)

	\bibitem{OEIS_A248044}
	OEIS Foundation:
	Sequence A248044.
	The On-Line Encyclopedia of Integer Sequences
	(accessed March~2026).
	\url{https://oeis.org/A248044}

	\bibitem{CortiA248044}
	Corti, C.:
	Computational verification of Sun's Conjecture~4.3(i)
	and exceptional structure of sequence~A248044.
	Preprint, in preparation (2026); data and code:
	\url{https://github.com/carcorti/A248044},
	\url{https://doi.org/10.5281/zenodo.19182594}

	\bibitem{GithubRepo}
	Corti, C.:
	A247975: computational study of Sun's Conjecture 4.1(i).
	GitHub repository and Zenodo archive (2026).
	\url{https://github.com/carcorti/A247975}.
	\url{https://doi.org/10.5281/zenodo.18920371}

\end{thebibliography}

\end{document}
